\section*{Chapitre \ref{ch:g7:chemical-reactions} --- Les réactions chimiques~: réactifs et produits}

\begin{solution}{exo:g7:chemical-reactions:1}
Physiques~: la glace qui fond, le sucre qui \omterm{def:g2:mixing-and-dissolving:dissolve}{se dissout}. Chimiques~:
l'allumette qui \omterm{def:g4:what-a-fire-needs:burning}{brûle}, le pain qui grille, le clou qui rouille.
\end{solution}

\begin{solution}{exo:g7:chemical-reactions:2}
Les \omterm{def:g7:chemical-reactions:reactant}{réactifs} sont les espèces consommées~; les \omterm{def:g7:chemical-reactions:reactant}{produits} sont les
nouvelles espèces formées.
\end{solution}

\begin{solution}{exo:g7:chemical-reactions:3}
carbone $+$ dioxygène $\longrightarrow$ dioxyde de carbone.
\end{solution}

\begin{solution}{exo:g7:chemical-reactions:4}
\omterm{def:g7:chemical-reactions:reactant}{Réactifs}~: le fer et le dioxygène. \omterm{def:g7:chemical-reactions:reactant}{Produit}~: l'oxyde de fer.
\end{solution}

\begin{solution}{exo:g7:chemical-reactions:5}
La buée révèle l'eau~; l'eau de chaux troublée révèle le dioxyde de
carbone.
\end{solution}

\begin{solution}{exo:g7:chemical-reactions:6}
zinc $+$ acide chlorhydrique $\longrightarrow$ chlorure de zinc $+$
dihydrogène.
\end{solution}

\begin{solution}{exo:g7:chemical-reactions:7}
Avant~: 4 (dans la \omterm{def:g7:atoms-and-molecules:molecule}{molécule} de méthane). Après~: $2 + 2 = 4$ (deux dans
chaque \omterm{def:g7:atoms-and-molecules:molecule}{molécule} d'eau).
\end{solution}

\begin{solution}{exo:g7:chemical-reactions:8}
Aucune espèce nouvelle n'apparaît~: le sel est toujours là, réparti dans
l'eau, et on le récupère en faisant évaporer l'eau. C'est une
\omterm{def:g7:chemical-reactions:physical-transformation}{transformation physique}.
\end{solution}

\begin{solution}{exo:g7:chemical-reactions:9}
La plus grande partie du bois a réagi avec le dioxygène pour donner des
\omterm{def:g7:chemical-reactions:reactant}{produits} qui sont des gaz, le dioxyde de carbone et la vapeur d'eau, qui
s'en vont dans l'\omterm{def:g4:air-a-mixture-of-gases:air}{air}. Seule la cendre reste.
\end{solution}

\begin{solution}{exo:g7:chemical-reactions:10}
dihydrogène $+$ dioxygène $\longrightarrow$ eau. Le \omterm{def:g7:chemical-reactions:reactant}{produit} bleuit le
sulfate de cuivre anhydre.
\end{solution}

\begin{solution}{exo:g7:chemical-reactions:11}
Avant~: $2 \times 2 = 4$ \omterm{def:g7:atoms-and-molecules:atom}{atomes} d'hydrogène et 2 \omterm{def:g7:atoms-and-molecules:atom}{atomes} d'oxygène.
Après~: deux \omterm{def:g7:atoms-and-molecules:molecule}{molécules} d'eau contiennent $2 \times 2 = 4$ \omterm{def:g7:atoms-and-molecules:atom}{atomes}
d'hydrogène et $2 \times 1 =
2$ \omterm{def:g7:atoms-and-molecules:atom}{atomes} d'oxygène. Les mêmes \omterm{def:g7:atoms-and-molecules:atom}{atomes}, en mêmes nombres.
\end{solution}

\begin{solution}{exo:g7:chemical-reactions:12}
Un \omterm{def:g7:atoms-and-molecules:atom}{atome} d'oxygène ne peut pas venir de nulle part~: le dessin a 2
\omterm{def:g7:atoms-and-molecules:atom}{atomes} d'oxygène avant et 3 après. Dessin correct~: un \omterm{def:g7:atoms-and-molecules:atom}{atome} de carbone
$+$ une \omterm{def:g7:atoms-and-molecules:molecule}{molécule} de dioxygène $\longrightarrow$ une \omterm{def:g7:atoms-and-molecules:molecule}{molécule} de dioxyde
de carbone, sans rien en trop.
\end{solution}

\begin{solution}{pb:g7:chemical-reactions:1}
\textbf{1.} Le méthane et le dioxygène.

\textbf{2.} L'eau.

\textbf{3.} Le dioxyde de carbone.

\textbf{4.} Non~: il n'est pas consommé, et il ne s'écrit pas dans
l'\omterm{def:g7:chemical-reactions:word-equation}{équation en mots}.

\textbf{5.} méthane $+$ dioxygène $\longrightarrow$ dioxyde de carbone
$+$ eau.

\textbf{6.} Chimique~: deux espèces disparaissent (le méthane, le
dioxygène) et deux espèces nouvelles apparaissent, mises en évidence par
des tests~; de plus, de la chaleur et de la lumière sont dégagées.

\textbf{7.} \ce{CH4}, \ce{O2}, \ce{CO2}, \ce{H2O}.

\textbf{8.} 1 carbone, 4 hydrogènes, $2 \times 2 = 4$ oxygènes.

\textbf{9.} 1 carbone~; $2 \times 2 = 4$ hydrogènes~; $2 + 2 \times 1 = 4$
oxygènes. Les mêmes \omterm{def:g7:atoms-and-molecules:atom}{atomes} en mêmes nombres~: ils ont seulement été
réarrangés.

\textbf{10.} $10 \times 2 = 20$ \omterm{def:g7:atoms-and-molecules:molecule}{molécules} de dioxygène.

\textbf{11.} 10 \omterm{def:g7:atoms-and-molecules:molecule}{molécules} de dioxyde de carbone.

\textbf{12.} $10 \times 2 = 20$ \omterm{def:g7:atoms-and-molecules:molecule}{molécules} d'eau.
\end{solution}
